\capitulo{v1c04}

%% =====================================================================
\section{Razão: comparar dividindo}
%% =====================================================================

Quantos candidatos disputam cada vaga? Quantos quilômetros um carro roda com
um litro de combustível? Qual time teve o melhor aproveitamento no campeonato?
Todas essas perguntas têm a mesma resposta matemática: uma \textbf{razão}.
Razão é a maneira de comparar duas quantidades por meio de uma
\emph{divisão} — e é uma das ideias mais cobradas no ENEM, quase sempre
disfarçada de ``taxa'', ``índice'', ``densidade'', ``rendimento'',
``aproveitamento'' ou ``proporção''.

Por que dividir, e não subtrair? Veja: Ana acertou 18 das 20 questões de um
simulado; Bruno acertou 40 das 50 questões de outro. Bruno acertou mais
questões ($40 > 18$), mas quem teve o melhor desempenho? A diferença não
responde a isso. A razão responde:
\[
  \text{Ana: } \frac{18}{20} = 0{,}90 = 90\%
  \qquad\qquad
  \text{Bruno: } \frac{40}{50} = 0{,}80 = 80\%.
\]
Ana acertou 9 de cada 10 questões; Bruno, 8 de cada 10. Quando os
``tamanhos'' das situações são diferentes, só a razão permite uma comparação
justa.

\begin{definicao}[Razão]
Dados dois números $a$ e $b$, com $b \neq 0$, a \textbf{razão de $a$ para $b$}
é o quociente
\[
  \frac{a}{b} \qquad\text{ou}\qquad a : b \qquad (\text{lê-se ``$a$ está para $b$''}).
\]
O número $a$ é o \textbf{antecedente} e $b$ é o \textbf{consequente}. A razão
$\frac{b}{a}$ (com $a \neq 0$) é a \textbf{razão inversa} de $\frac{a}{b}$.
\end{definicao}

A ordem importa! As expressões ``razão entre $A$ e $B$'', ``razão de $A$ para
$B$'', ``$A$ em relação a $B$'', ``$A$ por $B$'' e ``$A$ para cada $B$''
significam todas $\dfrac{A}{B}$: o que vem primeiro no texto vai para o
numerador.

\subsection{Razão entre grandezas de mesma natureza}

Quando as duas quantidades são da mesma espécie (comprimento com comprimento,
pessoas com pessoas, reais com reais), elas precisam estar na
\textbf{mesma unidade}. Aí as unidades se cancelam e a razão é um
\textbf{número puro}, que pode ser escrito como fração, decimal, porcentagem
ou na forma ``$1 : n$''.

Por exemplo, a razão entre a altura da tela de um celular (15~cm) e a altura de
um telão de shows (6~m) \emph{não} é $15/6$. Primeiro convertemos: $6~\text{m} = 600~\text{cm}$. Então
\[
  \frac{15~\text{cm}}{600~\text{cm}} = \frac{1}{40} = 1 : 40,
\]
ou seja, o telão é 40 vezes mais alto que a tela do celular.

\subsection{Comparando razões}

Para descobrir qual de duas ou mais razões é a maior, há três caminhos:
\begin{enumerate}[leftmargin=1.6em,itemsep=2pt]
  \item \textbf{Transformar em decimal ou porcentagem}: $\frac{3}{8} = 0{,}375$ e
        $\frac{5}{12} \approx 0{,}417$; logo $\frac{5}{12}$ é maior.
  \item \textbf{Igualar os consequentes} (denominadores): $\frac{3}{8} = \frac{9}{24}$
        e $\frac{5}{12} = \frac{10}{24}$; logo $\frac{5}{12}$ é maior.
  \item \textbf{Produto cruzado}: para $b, d > 0$, $\frac{a}{b} > \frac{c}{d}$ exatamente
        quando $a \cdot d > b \cdot c$. Como $3 \cdot 12 = 36 < 8 \cdot 5 = 40$, temos
        $\frac{3}{8} < \frac{5}{12}$.
\end{enumerate}

\begin{exemplo}[Aproveitamento nos lances livres]
Em um torneio de basquete, três jogadoras tiveram o seguinte desempenho nos
lances livres: Clara converteu 18 de 25 arremessos; Débora, 21 de 30; e Elisa,
26 de 40. Qual delas teve o melhor aproveitamento? E qual converteu mais
lances livres?
\solucao
O aproveitamento é a razão entre os lances convertidos e os lances tentados:
\[
  \text{Clara: } \frac{18}{25} = 0{,}72 \qquad
  \text{Débora: } \frac{21}{30} = 0{,}70 \qquad
  \text{Elisa: } \frac{26}{40} = 0{,}65.
\]
Clara teve o melhor aproveitamento (72\%), embora Elisa tenha convertido mais
lances (26). As duas perguntas têm respostas \emph{diferentes} — e o ENEM adora
colocar a resposta da pergunta errada entre as alternativas. Conferindo Clara
e Débora pelo produto cruzado: $18 \cdot 30 = 540 > 25 \cdot 21 = 525$.
\end{exemplo}

\subsection{Parte-parte e parte-todo}

Considere a frase ``em uma turma, há 3 meninas para cada 2 meninos''. Ela
compara uma \textbf{parte} com outra \textbf{parte}: meninas $:$ meninos $= 3 : 2$.
Mas que fração da turma é formada por meninas? Basta juntar as partes: em cada
grupo de $3 + 2 = 5$ estudantes, 3 são meninas.

\begin{center}
\begin{tikzpicture}[x=1cm,y=1cm]
  \foreach \i in {0,1,2}{
    \node[fill=ebookPrim,text=white,rounded corners=2pt,minimum size=9mm,font=\large]
      at (\i*1.05,0) {\faFemale};}
  \foreach \i in {3,4}{
    \node[fill=ebookAcc,text=white,rounded corners=2pt,minimum size=9mm,font=\large]
      at (\i*1.05,0) {\faMale};}
  \draw[decorate,decoration={brace,amplitude=5pt,mirror},thick,ebookPrim]
    (-0.45,-0.6) -- (2.55,-0.6) node[midway,below=6pt,font=\small\sffamily]{meninas: 3 partes};
  \draw[decorate,decoration={brace,amplitude=5pt,mirror},thick,ebookAcc!80!black]
    (2.7,-0.6) -- (4.65,-0.6) node[midway,below=6pt,font=\small\sffamily]{meninos: 2 partes};
  \draw[decorate,decoration={brace,amplitude=5pt},thick,ebookMuted]
    (-0.45,0.6) -- (4.65,0.6) node[midway,above=6pt,font=\small\sffamily]{turma: $3+2=5$ partes};
  \node[anchor=west,font=\small,align=left] at (5.4,0.1)
    {parte-parte: $\dfrac{\text{meninas}}{\text{meninos}} = \dfrac{3}{2}$\\[8pt]
     parte-todo: $\dfrac{\text{meninas}}{\text{turma}} = \dfrac{3}{5} = 60\%$};
\end{tikzpicture}
\end{center}

\begin{formula}[Da razão parte-parte à fração do todo]
Se $A : B = m : n$, então
$\displaystyle \frac{A}{A+B} = \frac{m}{m+n}$ \quad e \quad
$\displaystyle \frac{B}{A+B} = \frac{n}{m+n}$.
\end{formula}

\begin{exemplo}[Aprovados e reprovados]
Em um concurso, a razão entre o número de candidatos aprovados e o de
reprovados foi de $2 : 7$. Sabendo que 1\,260 candidatos fizeram a prova,
quantos foram aprovados?
\solucao
A razão $2 : 7$ é parte-parte. Juntando as partes, o total corresponde a
$2 + 7 = 9$ partes iguais. Cada parte vale $1\,260 \div 9 = 140$ candidatos.
Logo, foram aprovados $2 \cdot 140 = \mathbf{280}$ candidatos (e reprovados
$7 \cdot 140 = 980$; confira: $280 + 980 = 1\,260$).

\textbf{Cuidado:} calcular $\frac{2}{7}$ de $1\,260$ (que dá 360) é tratar a razão
parte-parte como se fosse parte-todo — um distrator clássico.
\end{exemplo}

%% =====================================================================
\section{Taxas: razões com unidade}
%% =====================================================================

Quando as duas quantidades comparadas são de \emph{naturezas diferentes}, a
razão carrega uma unidade, e a lemos com a palavra ``por'': quilômetros
\emph{por} hora, habitantes \emph{por} quilômetro quadrado, miligramas
\emph{por} mililitro. Essas razões costumam ser chamadas de \textbf{taxas}.

A grande vantagem: \textbf{a unidade diz qual é a conta}. Se a resposta deve
sair em km/h, é porque você vai dividir uma distância em km por um tempo em
horas. Se sair em L/100\,km, divide litros por ``centenas de quilômetros''.

\begin{formula}[As taxas mais cobradas no ENEM]
\small
\begin{tabularx}{\linewidth}{>{\raggedright\arraybackslash}p{3.4cm}>{\raggedright\arraybackslash}X>{\raggedright\arraybackslash}p{2.9cm}}
\toprule
\textbf{Taxa} & \textbf{Razão} & \textbf{Unidades usuais}\\
\midrule
Velocidade média & distância $\div$ tempo & km/h, m/s\\
Densidade (de um material) & massa $\div$ volume & g/cm³, kg/L, kg/m³\\
Densidade demográfica & número de habitantes $\div$ área & hab/km²\\
Rendimento de combustível & distância $\div$ volume de combustível & km/L\\
Consumo de combustível & volume de combustível $\div$ distância & L/100\,km\\
Vazão & volume $\div$ tempo & L/min, m³/h\\
Concentração & massa de soluto $\div$ volume da solução & g/L, mg/mL\\
Produtividade & produção $\div$ área (ou $\div$ trabalhador) & t/ha\\
Renda \emph{per capita} & renda total $\div$ número de pessoas & R\$/pessoa\\
\bottomrule
\end{tabularx}
\end{formula}

\subsection{Velocidade média}

A velocidade média é a razão entre a distância percorrida e o tempo total gasto
para percorrê-la — \emph{incluindo} paradas, se elas fizerem parte do intervalo
considerado.

\begin{formula}[Velocidade média]
$\displaystyle v_m = \frac{\Delta s}{\Delta t}
 \qquad\qquad 1~\text{m/s} = \frac{3\,600~\text{m}}{1~\text{h}} = 3{,}6~\text{km/h}$
\end{formula}

No gráfico da posição em função do tempo, a velocidade média é a
\emph{inclinação} do segmento que liga o ponto de partida ao de chegada. O
gráfico abaixo representa a viagem do próximo exemplo.

\begin{center}
\begin{tikzpicture}
\begin{axis}[width=10.5cm,height=6.2cm,xmin=0,xmax=3.3,ymin=0,ymax=230,
  axis lines=left,xlabel={tempo (h)},ylabel={posição (km)},
  xtick={0,0.5,1,1.5,2,2.5,3},ytick={0,40,80,120,160,200},
  xticklabel style={/pgf/number format/use comma},
  grid=major,grid style={ebookLine!60},
  label style={font=\small\sffamily},tick label style={font=\scriptsize\sffamily}]
  \addplot[very thick,ebookPrim,mark=*,mark size=1.6pt] coordinates {(0,0) (1.5,120) (2,120) (3,200)};
  \addplot[thick,dashed,ebookAcc] coordinates {(0,0) (3,200)};
  \node[font=\scriptsize\sffamily,fill=white,inner sep=1.5pt,anchor=south east] at (axis cs:1.45,125) {80~km/h};
  \node[font=\scriptsize\sffamily,fill=white,inner sep=1.5pt,anchor=south] at (axis cs:1.75,122) {parada};
  \node[font=\scriptsize\sffamily,fill=white,inner sep=1.5pt,anchor=north west] at (axis cs:2.55,158) {80~km/h};
  \node[font=\scriptsize\sffamily,text=ebookAcc!70!black,fill=white,inner sep=1.5pt,anchor=west] at (axis cs:1.55,75)
    {$v_m = \frac{200}{3} \approx 66{,}7$~km/h};
\end{axis}
\end{tikzpicture}
\end{center}

\begin{exemplo}[Viagem com parada]
Uma família percorreu 120~km em 1~h~30~min, parou 30~min para almoçar e,
depois, percorreu os 80~km restantes em 1~h. Qual foi a velocidade média na
viagem inteira?
\solucao
A distância total é $\Delta s = 120 + 80 = 200$~km. O tempo total, incluindo a
parada, é $\Delta t = 1{,}5 + 0{,}5 + 1 = 3$~h (lembre: 30~min $= 0{,}5$~h). Então
\[
  v_m = \frac{200~\text{km}}{3~\text{h}} \approx 66{,}7~\text{km/h}.
\]
Repare que, nos dois trechos em movimento, a família andou a 80~km/h
($120 \div 1{,}5$ e $80 \div 1$), mas a parada ``puxou'' a velocidade média para
baixo. Velocidade média não é a média das velocidades!
\end{exemplo}

Outro caso clássico: ir a uma cidade a 60~km de distância a 60~km/h (1~h) e
voltar a 30~km/h (2~h). A velocidade média é $\frac{120~\text{km}}{3~\text{h}} = 40$~km/h,
e não $\frac{60+30}{2} = 45$~km/h. Como o carro passa \emph{mais tempo} no
trecho lento, esse trecho ``pesa mais'' no resultado.

\subsection{Densidade}

A \textbf{densidade} de um material é a razão entre sua massa e seu volume,
$d = \dfrac{m}{V}$. Ela diz ``quanto de massa cabe em cada unidade de volume''.
Algumas equivalências muito úteis:
\[
  1~\text{g/cm}^3 = 1~\text{kg/L} = 1\,000~\text{kg/m}^3 .
\]
A água tem densidade de aproximadamente 1~kg/L; a gasolina, cerca de 0,75~kg/L;
o óleo de cozinha, cerca de 0,9~kg/L (por isso flutua na água). A mesma ideia
aparece em outros contextos:
\begin{itemize}[leftmargin=1.4em,itemsep=1pt]
  \item \textbf{densidade demográfica} $=\dfrac{\text{habitantes}}{\text{área}}$ (hab/km²);
  \item \textbf{densidade superficial} $=\dfrac{\text{massa}}{\text{área}}$ (por exemplo, papel de 75~g/m²).
\end{itemize}

\begin{exemplo}[Carga de um caminhão-tanque]
Um caminhão-tanque, cuja massa vazio é de 16~t, transporta \num{25000}~L de
etanol, que tem densidade de 0,79~kg/L. Nessa rodovia, o peso bruto total
máximo permitido (caminhão mais carga) é de 36~t. O caminhão está dentro do
limite? Qual é o maior volume de etanol que ele poderia transportar?
\solucao
Da densidade, $m = d \cdot V$:
\[
  m = 0{,}79~\tfrac{\text{kg}}{\text{L}} \cdot 25\,000~\text{L} = 19\,750~\text{kg} = 19{,}75~\text{t}.
\]
Total: $16 + 19{,}75 = 35{,}75$~t $\leq 36$~t. Está dentro do limite.
A carga máxima é $36 - 16 = 20$~t $= 20\,000$~kg, o que corresponde a
$V = \dfrac{m}{d} = \dfrac{20\,000}{0{,}79} \approx 25\,316$~L.
Note como as unidades guiaram a conta: kg/L vezes L dá kg; kg dividido por
kg/L dá L.
\end{exemplo}

\subsection{Consumo, vazão e taxas ``por mil''}

\textbf{Rendimento e consumo de combustível.} Há duas formas de informar a
eficiência de um veículo, e elas são \emph{inversas} uma da outra:
o rendimento $r$, em km/L (quantos quilômetros por litro), e o consumo $c$, em
L/100\,km (quantos litros para rodar 100~km). Como 100~km exigem
$100 \div r$ litros,
\[
  c = \frac{100}{r} \qquad\text{e}\qquad r = \frac{100}{c}.
\]
Assim, 12,5~km/L equivale a 8~L/100\,km. \textbf{Maior rendimento significa
menor consumo.} O gráfico mostra essa relação inversa.

\begin{center}
\begin{tikzpicture}
\begin{axis}[width=9.5cm,height=5.6cm,xmin=4,xmax=26,ymin=0,ymax=26,
  axis lines=left,xlabel={rendimento $r$ (km/L)},ylabel={consumo $c$ (L/100 km)},
  xtick={5,10,15,20,25},ytick={0,5,10,15,20,25},
  grid=major,grid style={ebookLine!60},
  label style={font=\small\sffamily},tick label style={font=\scriptsize\sffamily}]
  \addplot[very thick,ebookPrim,domain=4.2:25,samples=80] {100/x};
  \addplot[only marks,mark=*,mark size=2pt,ebookAcc] coordinates {(5,20) (10,10) (12.5,8) (20,5)};
  \node[font=\scriptsize\sffamily,anchor=south west] at (axis cs:10,10) {(10; 10)};
  \node[font=\scriptsize\sffamily,anchor=south west] at (axis cs:12.5,8) {(12,5; 8)};
  \node[font=\scriptsize\sffamily,anchor=south west] at (axis cs:20,5) {(20; 5)};
  \node[font=\scriptsize\sffamily,anchor=west] at (axis cs:5.3,20) {(5; 20)};
\end{axis}
\end{tikzpicture}
\end{center}

\textbf{Vazão.} É a razão entre o volume que passa (ou que entra em um
reservatório) e o tempo: $Q = \dfrac{V}{t}$. Um chuveiro com vazão de
12~L/min, ligado durante um banho de 15~min, gasta $12 \cdot 15 = 180$~L.

\textbf{Taxas ``por mil'' e ``por 100 mil''.} Para comparar cidades com
populações muito diferentes, órgãos de saúde e de segurança usam índices como
``casos por 100 mil habitantes''. Uma cidade com 45 casos em \num{300000}
habitantes tem taxa de
\[
  \frac{45}{300\,000} = \frac{15}{100\,000} \;\Longrightarrow\; 15~\text{casos por 100 mil habitantes.}
\]

\begin{contexto}[Por que os jornais falam em ``por 100 mil habitantes''?]
Uma capital com 3 milhões de habitantes e 600 casos de uma doença tem mais
casos do que uma cidade de 50 mil habitantes com 40 casos. Mas o risco é maior
na cidade pequena: $600 \div 30 = 20$ casos por 100 mil contra
$40 \div 0{,}5 = 80$ casos por 100 mil. Sem a razão, a comparação engana.
\end{contexto}

%% =====================================================================
\section{Proporção e suas propriedades}
%% =====================================================================

Quando duas razões são iguais, dizemos que elas formam uma \textbf{proporção}.
É o que acontece quando ampliamos uma receita, quando mantemos o ``sabor'' de
um refresco ou quando duas cidades têm a mesma densidade demográfica.

\begin{definicao}[Proporção]
Os números $a, b, c, d$ (com $b, d \neq 0$) formam uma \textbf{proporção},
nessa ordem, quando
\[
  \frac{a}{b} = \frac{c}{d} \qquad (\text{lê-se ``$a$ está para $b$ assim como $c$ está para $d$''}).
\]
Os termos $a$ e $d$ são os \textbf{extremos}; $b$ e $c$ são os \textbf{meios}.
\end{definicao}

\begin{formula}[Propriedade fundamental]
$\displaystyle \frac{a}{b} = \frac{c}{d} \iff a \cdot d = b \cdot c$
\qquad (produto dos extremos $=$ produto dos meios)
\end{formula}

É ela que permite achar um termo desconhecido (a ``quarta proporcional''): se
$\frac{x}{12} = \frac{5}{8}$, então $8x = 60$ e $x = 7{,}5$.

\subsection{Outras propriedades úteis}

\begin{formula}[Propriedades das proporções]
Se $\displaystyle \frac{a}{b} = \frac{c}{d}$, então:\\[6pt]
$\displaystyle \text{(P1)}\quad \frac{a+b}{b} = \frac{c+d}{d}
  \quad\text{e}\quad \frac{a-b}{b} = \frac{c-d}{d}$\\[8pt]
$\displaystyle \text{(P2)}\quad \frac{a+c}{b+d} = \frac{a-c}{b-d} = \frac{a}{b} = \frac{c}{d}
  \quad (\text{com } b+d \neq 0 \text{ e } b-d \neq 0)$\\[8pt]
Em geral, se $\dfrac{a}{x} = \dfrac{b}{y} = \dfrac{c}{z} = k$, então
$\dfrac{a+b+c}{x+y+z} = k$.
\end{formula}

A propriedade (P2) diz que, em uma sequência de razões iguais, ``somar os
numeradores e somar os denominadores'' preserva a razão. É o que justifica o
método das partes: se $\frac{a}{x} = \frac{b}{y} = k$, então $a = kx$ e $b = ky$.

\begin{exemplo}[Eleição no condomínio]
Na eleição para síndico de um condomínio, os votos dos candidatos A e B ficaram
na razão $5 : 3$, e A recebeu 84 votos a mais que B. Quantos votos teve cada
candidato?
\solucao
Escreva $A = 5k$ e $B = 3k$. A diferença é $A - B = 2k = 84$, logo $k = 42$.
Assim, $A = 5 \cdot 42 = \mathbf{210}$ votos e $B = 3 \cdot 42 = \mathbf{126}$ votos.

Pela propriedade (P2): $\dfrac{A}{5} = \dfrac{B}{3} = \dfrac{A - B}{5 - 3} = \dfrac{84}{2} = 42$.
É a mesma conta, agora justificada.
\end{exemplo}

\subsection{Razões de mistura}

Na construção civil, na cozinha, na farmácia e na agricultura, misturas são
descritas por razões como ``1 : 2 : 6''. Isso significa que a mistura é formada
por \textbf{partes iguais de volume (ou de massa)}: 1 parte do primeiro
componente, 2 do segundo e 6 do terceiro — ao todo, $1 + 2 + 6 = 9$ partes.

\begin{center}
\begin{tikzpicture}[x=1.15cm,y=1cm]
  \fill[ebookMuted!70] (0,0) rectangle (1,0.8);
  \fill[ebookLine] (1,0) rectangle (3,0.8);
  \fill[ebookAcc!55] (3,0) rectangle (9,0.8);
  \foreach \i in {0,...,9}{\draw[white,line width=1.2pt] (\i,0) -- (\i,0.8);}
  \draw[ebookText!60,line width=.5pt,rounded corners=1pt] (0,0) rectangle (9,0.8);
  \draw[decorate,decoration={brace,amplitude=4pt,mirror},thick] (0.02,-0.12) -- (0.98,-0.12)
    node[midway,below=5pt,font=\scriptsize\sffamily,align=center]{cimento\\1 parte};
  \draw[decorate,decoration={brace,amplitude=4pt,mirror},thick] (1.02,-0.12) -- (2.98,-0.12)
    node[midway,below=5pt,font=\scriptsize\sffamily,align=center]{cal\\2 partes};
  \draw[decorate,decoration={brace,amplitude=4pt,mirror},thick] (3.02,-0.12) -- (8.98,-0.12)
    node[midway,below=5pt,font=\scriptsize\sffamily,align=center]{areia\\6 partes};
  \draw[decorate,decoration={brace,amplitude=5pt},thick,ebookPrim] (0,0.92) -- (9,0.92)
    node[midway,above=6pt,font=\small\sffamily,text=ebookPrim]{argamassa: $1+2+6 = 9$ partes iguais};
\end{tikzpicture}
\end{center}

\begin{formula}[Método das partes]
valor de uma parte $=\dfrac{\text{total da mistura}}{\text{soma das partes}}$
\qquad;\qquad
componente $=$ (número de partes) $\times$ (valor de uma parte)
\end{formula}

\begin{exemplo}[Argamassa de assentamento]
Uma argamassa é preparada com cimento, cal e areia na razão $1 : 2 : 6$, em
volume. Para obter 4,5~m³ de argamassa, quantos metros cúbicos de cada material
são necessários?
\solucao
São $1 + 2 + 6 = 9$ partes. Cada parte vale $4{,}5 \div 9 = 0{,}5$~m³. Assim:
cimento $= 1 \cdot 0{,}5 = 0{,}5$~m³; cal $= 2 \cdot 0{,}5 = 1$~m³;
areia $= 6 \cdot 0{,}5 = 3$~m³. Conferência: $0{,}5 + 1 + 3 = 4{,}5$~m³.
\end{exemplo}

E quando juntamos \emph{duas misturas diferentes}? Aí não se somam razões:
somam-se as \textbf{quantidades} de cada componente.

\begin{exemplo}[Juntando dois refrescos]
A jarra A contém 2~L de um refresco preparado com xarope e água na razão
$1 : 4$; a jarra B contém 3~L de refresco na razão $1 : 2$. Despejando as duas
jarras em um recipiente maior, qual será a razão entre xarope e água da
mistura?
\solucao
Jarra A: $1 + 4 = 5$ partes de $2 \div 5 = 0{,}4$~L $\Rightarrow$ xarope 0,4~L e água 1,6~L.\\
Jarra B: $1 + 2 = 3$ partes de $3 \div 3 = 1$~L $\Rightarrow$ xarope 1~L e água 2~L.\\
Mistura: xarope $= 0{,}4 + 1 = 1{,}4$~L e água $= 1{,}6 + 2 = 3{,}6$~L. A razão é
\[
  \frac{1{,}4}{3{,}6} = \frac{14}{36} = \frac{7}{18}, \quad\text{ou seja, } 7 : 18.
\]
Somar as razões ``de cabeça'' — $(1+1) : (4+2) = 1 : 3$ — está \textbf{errado},
porque as jarras têm volumes diferentes. A propriedade (P2) só vale quando as
razões somadas são \emph{iguais}.
\end{exemplo}

%% =====================================================================
\section{Dicas e macetes}
%% =====================================================================

\begin{dica}[Leia a razão na ordem do texto]
``Razão entre $X$ e $Y$'', ``$X$ em relação a $Y$'', ``$X$ por $Y$'', ``$X$ para cada
$Y$'': em todos os casos, $X$ vai no numerador. Antes de fazer contas, escreva a
razão \emph{em palavras} — por exemplo, $\dfrac{\text{assentos vendidos}}{\text{total de assentos}}$.
As alternativas costumam trazer a razão invertida como distrator.
\end{dica}

\begin{dica}[Deixe a unidade guiar a conta]
Se a resposta é pedida em gotas por minuto, você vai dividir um número de gotas
por um número de minutos. Se é pedida em kg, e você tem kg/L e L, vai multiplicar
($\frac{\text{kg}}{\text{L}} \cdot \text{L} = \text{kg}$). Essa ``análise
dimensional'' evita a maioria dos erros de montagem.
\end{dica}

\begin{dica}[Só há relações? Atribua valores]
Quando o enunciado não dá números, apenas relações (``a área de $Q$ é três
quartos da área de $R$'', ``a massa de $A$ é 1,5 vez a de $B$''), invente valores
convenientes: área de $R$ igual a 4, população igual a 100\ldots{} Calcule as razões
com esses números e compare. O resultado não depende dos valores escolhidos.
\end{dica}

\begin{dica}[Gráfico de barras com eixo cortado]
A razão entre as \emph{alturas} de duas barras só é igual à razão entre os
\emph{valores} quando o eixo vertical começa em zero. Se o eixo começa em 20,
uma barra de valor 30 aparece com altura proporcional a $30 - 20 = 10$.
\end{dica}

\begin{macete}[Produto cruzado para comparar frações]
Para $b, d > 0$: $\dfrac{a}{b} > \dfrac{c}{d} \iff a \cdot d > b \cdot c$.
Exemplo: $\frac{7}{12}$ ou $\frac{11}{19}$? Como $7 \cdot 19 = 133 > 12 \cdot 11 = 132$,
temos $\frac{7}{12} > \frac{11}{19}$ — sem calcular nenhum decimal.
\end{macete}

\begin{macete}[km/h e m/s: o fator 3,6]
De m/s para km/h, multiplique por 3,6; de km/h para m/s, divida por 3,6.
Exemplos: $20~\text{m/s} = 72~\text{km/h}$; $\;90~\text{km/h} = 25~\text{m/s}$.
Lembre também: $1~\text{h}\,12~\text{min} = 1{,}2~\text{h}$, porque
$12~\text{min} = \frac{12}{60}~\text{h}$.
\end{macete}

\begin{macete}[Ida e volta: a média harmônica]
Se um percurso de ida e um de volta têm a \emph{mesma distância} e são feitos às
velocidades $v_1$ e $v_2$, a velocidade média total é
$\displaystyle v_m = \frac{2\,v_1 v_2}{v_1 + v_2}$.
Exemplo: 60~km/h e 30~km/h dão $\frac{2 \cdot 60 \cdot 30}{90} = 40$~km/h.
\end{macete}

\begin{atencao}[Parte-parte não é parte-todo]
Se a razão entre aprovados e reprovados é $2 : 7$, os aprovados são
$\frac{2}{9}$ do total, e não $\frac{2}{7}$. Sempre que precisar da fração do
todo, \textbf{some as partes} para obter o denominador.
\end{atencao}

\begin{atencao}[Velocidade média não é média das velocidades]
A velocidade média é sempre \emph{distância total} dividida por \emph{tempo
total}, incluindo paradas que estejam dentro do intervalo. A média aritmética
das velocidades dos trechos só coincide com ela quando os trechos duram o mesmo
\emph{tempo}.
\end{atencao}

\begin{atencao}[Unidades antes da divisão]
Em razões entre grandezas de mesma natureza, converta tudo para a mesma
unidade (cm com cm, minuto com minuto). E não confunda km/L com L/100\,km:
são grandezas inversas — quanto maior o rendimento, menor o consumo.
\end{atencao}

\begin{resumo}
\begin{itemize}[leftmargin=1.3em,itemsep=2pt]
  \item \textbf{Razão} de $a$ para $b$: $\frac{a}{b}$ (o primeiro citado vai no numerador).
        Mesma natureza $\Rightarrow$ mesma unidade e número puro.
  \item \textbf{Comparar razões}: decimal/porcentagem, denominador comum ou produto cruzado.
  \item \textbf{Parte-parte} $m : n$ $\Rightarrow$ fração do todo $\frac{m}{m+n}$.
  \item \textbf{Taxas}: $v_m = \frac{\Delta s}{\Delta t}$; $d = \frac{m}{V}$;
        densidade demográfica $=\frac{\text{hab}}{\text{área}}$; vazão $=\frac{V}{t}$;
        consumo $c = \frac{100}{r}$ (L/100\,km $\leftrightarrow$ km/L).
  \item $1~\text{m/s} = 3{,}6~\text{km/h}$; $\;1~\text{g/cm}^3 = 1~\text{kg/L} = 1\,000~\text{kg/m}^3$.
  \item \textbf{Proporção}: $\frac{a}{b} = \frac{c}{d} \iff ad = bc$; razões iguais
        $\Rightarrow$ $\frac{a+c}{b+d}$ também é igual a elas.
  \item \textbf{Misturas}: valor de uma parte $=$ total $\div$ soma das partes; ao juntar
        misturas, some \emph{quantidades}, nunca razões.
\end{itemize}
\end{resumo}

%% =====================================================================
\secaoenem
%% =====================================================================

\questaoenem{2014-154}
\begin{resolucao}{D}
\passo{1. Entendendo o problema.} O próprio texto define o desempenho: é a razão
entre o número de vezes em que o jogador derrubou todos os pinos e o número de
jogadas. Precisamos calcular essa razão para os cinco jogadores e escolher a maior.

\passo{2. Calculando as razões.}
\[
\text{I: } \tfrac{50}{85} \approx 0{,}588 \quad
\text{II: } \tfrac{40}{65} \approx 0{,}615 \quad
\text{III: } \tfrac{20}{65} \approx 0{,}308 \quad
\text{IV: } \tfrac{30}{40} = 0{,}75 \quad
\text{V: } \tfrac{48}{90} \approx 0{,}533
\]

\passo{3. Conclusão.} O maior desempenho é o do jogador IV: ele derrubou todos os
pinos em 75\% das jogadas. Repare que IV foi o que \emph{menos jogou} e um dos que
\emph{menos} derrubaram todos os pinos em números absolutos — é exatamente por isso
que se usa a razão. Item fácil ($b = -0{,}36$): a razão está definida no
enunciado e basta aplicá-la.
\resposta{alternativa D — jogador IV.}
\begin{distratores}
\alt{A} Quem compara apenas os números absolutos escolhe o jogador I, que derrubou
todos os pinos mais vezes (50). Mas ele também jogou muito (85 vezes): $50/85 \approx 0{,}59$.
\alt{B} O jogador II tem a segunda maior razão ($\approx 0{,}62$). Quem estima ``no
olho'' em vez de calcular, ou compara apenas II e III (mesmo número de jogadas), pode
parar aqui.
\alt{C} Quem inverte a razão (jogadas $\div$ acertos) encontra o maior valor para o
jogador III: $65/20 = 3{,}25$. Mas essa razão mede quantas jogadas ele precisa para
cada acerto — quanto maior, \emph{pior}.
\alt{E} O jogador V tem 48 acertos e o maior número de jogadas (90), o que pode
impressionar; porém sua razão é só $\approx 0{,}53$, a segunda menor.
\end{distratores}
\end{resolucao}

\questaoenem{2025-176}
\begin{resolucao}{C}
\passo{1. Entendendo o problema.} O jogador participou do segundo tempo (45~min) e da
prorrogação (30~min). No segundo tempo, percorreu 4,5~km; na prorrogação, manteve a
\emph{mesma velocidade média}. Queremos a distância total durante a participação dele.

\passo{2. Velocidade média no segundo tempo.}
\[
v_m = \frac{4{,}5~\text{km}}{45~\text{min}} = 0{,}1~\text{km/min} \quad(= 6~\text{km/h}).
\]

\passo{3. Distância na prorrogação e total.} Em 30~min, à mesma taxa:
$0{,}1 \cdot 30 = 3$~km. Total: $4{,}5 + 3 = 7{,}5$~km.

Atalho: como a velocidade é a mesma, distância e tempo são proporcionais; ele jogou
$45 + 30 = 75$~min no total, ou seja, $\frac{75}{45} = \frac{5}{3}$ do tempo do segundo
tempo. Logo, percorreu $\frac{5}{3} \cdot 4{,}5 = 7{,}5$~km. Item de dificuldade média
($b = 1{,}06$): exige notar que o jogador \emph{não} jogou o primeiro tempo.
\resposta{alternativa C — 7,5~km.}
\begin{distratores}
\alt{A} 4,5~km é apenas a distância do segundo tempo; esquece a prorrogação.
\alt{B} 6,0 é a velocidade média em km/h ($4{,}5 \div 0{,}75$~h), não a distância.
Confundir a taxa com a grandeza pedida é um erro comum.
\alt{D} 9,0~km supõe que a prorrogação durou tanto quanto o segundo tempo
($4{,}5 + 4{,}5$), ignorando que ela tem só 30~min.
\alt{E} 12,0~km corresponde a 120~min à taxa de 0,1~km/min — como se o jogador tivesse
participado da partida inteira, inclusive do primeiro tempo.
\end{distratores}
\end{resolucao}

\questaoenem{2013-159}
\begin{resolucao}{B}
\passo{1. Entendendo a proporção.} ``1 parte de cimento, 4 de areia e 2 de brita''
significa que o concreto é formado por $1 + 4 + 2 = 7$ partes iguais de volume.

\passo{2. Valor de cada parte.} $14~\text{m}^3 \div 7 = 2~\text{m}^3$ por parte.

\passo{3. Cimento.} O cimento corresponde a 1 parte: $1 \cdot 2 = 2$~m³.
(Areia: $4 \cdot 2 = 8$~m³; brita: $2 \cdot 2 = 4$~m³; total: 14~m³.)
Item de dificuldade média ($b = 1{,}15$): é a aplicação direta do método das partes,
mas exige transformar a razão parte-parte em fração do todo ($\frac{1}{7}$).
\resposta{alternativa B — 2,00~m³.}
\begin{distratores}
\alt{A} 1,75 é $14 \div 8$: surge de quem conta errado o total de partes (por
exemplo, multiplicando $4 \cdot 2$ em vez de somar $1 + 4 + 2$).
\alt{C} 2,33 é $14 \div 6$: o aluno divide pelo número de partes de areia e brita
($4 + 2$), esquecendo de incluir a própria parte de cimento no total.
\alt{D} 4,00~m³ é o volume de \emph{brita} (2 partes), não de cimento.
\alt{E} 8,00~m³ é o volume de \emph{areia} (4 partes).
\end{distratores}
\end{resolucao}

\questaoenem{2024-166}
\begin{resolucao}{E}
\passo{1. Escrevendo as densidades.} Sejam $H$ o número de habitantes de $R$ e $A$ a
área de $R$. Então $d(R) = \dfrac{H}{A}$. Para $Q$: habitantes $=\dfrac{H}{2}$ e área
$=\dfrac{3A}{4}$.

\passo{2. Calculando $d(Q)$.}
\[
d(Q) = \frac{H/2}{3A/4} = \frac{H}{2} \cdot \frac{4}{3A} = \frac{4}{6} \cdot \frac{H}{A} = \frac{2}{3}\,d(R).
\]

\passo{3. Conferindo com números (macete de atribuir valores).} Suponha que $R$ tem
área 4~km² e 120 habitantes: $d(R) = 30$. Então $Q$ tem área 3~km² e 60 habitantes:
$d(Q) = 20 = \frac{2}{3} \cdot 30$. Item médio ($b = 1{,}74$): a dificuldade está em
lidar com uma fração \emph{dividida} por outra fração — o número de habitantes cai à
metade, mas a área também diminui, o que ``compensa'' parcialmente.
\resposta{alternativa E — $d(Q) = \frac{2}{3}\,d(R)$.}
\begin{distratores}
\alt{A} $\frac{1}{4}$ vem de subtrair as frações ($\frac{3}{4} - \frac{1}{2}$), operação
que não tem relação com a definição de densidade.
\alt{B} $\frac{1}{2}$ considera apenas a redução do número de habitantes e esquece que a
área de $Q$ também é menor.
\alt{C} $\frac{3}{4}$ considera apenas a área, e ainda como se a densidade fosse
\emph{diretamente} proporcional a ela (na verdade, a área está no denominador).
\alt{D} $\frac{3}{2}$ é a razão invertida: $\frac{3/4}{1/2}$, área sobre habitantes.
\end{distratores}
\end{resolucao}

\questaoenem{2020-147}
\begin{resolucao}{C}
\passo{1. Interpretando a tela do tipo A.} A informação ``$X/100$'' significa que o
veículo gastou $X$ litros para percorrer 100~quilômetros: é um consumo em L/100\,km.

\passo{2. O que o tipo B informa.} O tipo B mostra quantos quilômetros o veículo
percorre com \emph{um} litro (rendimento em km/L). É a razão inversa:
\[
\frac{X~\text{litros}}{100~\text{km}} \;\Longrightarrow\; \frac{100~\text{km}}{X~\text{litros}} = \frac{100}{X}~\text{km/L}.
\]
Conferindo com um número: se $X = 8$ (8~L a cada 100~km), o carro roda
$100 \div 8 = 12{,}5$~km com 1~L. De fato, $\frac{100}{8} = 12{,}5$.

\passo{3. Por que é difícil.} Item difícil ($b = 2{,}78$): é uma questão de
\emph{taxa} em linguagem algébrica, sem números. A notação ``$X/100$'' na tela
induz o aluno a repeti-la ou a operá-la sem pensar no significado de cada termo.
Testar com um valor concreto, como fizemos, desfaz a confusão.
\resposta{alternativa C — $\dfrac{100}{X}$.}
\begin{distratores}
\alt{A} $X \cdot 100$ multiplica as grandezas em vez de inverter a razão; com $X = 8$,
daria 800~km/L, um absurdo.
\alt{B} $\frac{X}{100}$ apenas repete a informação do tipo A (litros por quilômetro),
que é o inverso do que se pede.
\alt{D} $\frac{1}{X}$ inverte, mas esquece os 100~km: seria correto apenas se $X$ fosse o
gasto para 1~km.
\alt{E} $1 \cdot X = X$ trata o consumo em L/100\,km como se já fosse o rendimento em km/L.
\end{distratores}
\end{resolucao}

\questaoenem{2017-165}
\begin{resolucao}{E}
\passo{1. Alturas no Gráfico 1.} O eixo começa em 0, então as alturas das colunas
são os próprios valores: A $= 70$ e B $= 30$. Razão B/A:
$\frac{30}{70} = \frac{3}{7}$.

\passo{2. Alturas no Gráfico 2.} O eixo foi cortado e começa em 20. A altura
\emph{visível} de cada coluna é o valor menos 20: A $\to 70 - 20 = 50$ e
B $\to 30 - 20 = 10$. Razão B/A: $\frac{10}{50} = \frac{1}{5}$.

\passo{3. Diferença.}
\[
\frac{3}{7} - \frac{1}{5} = \frac{15}{35} - \frac{7}{35} = \frac{8}{35}.
\]
No Gráfico 2, a coluna B parece ter $\frac{1}{5}$ da altura de A, quando deveria ter
$\frac{3}{7}$ (quase metade) — daí a crítica dos leitores. Item muito difícil
($b = 3{,}63$): o aluno precisa perceber que a pergunta é sobre \emph{alturas
desenhadas}, não sobre os percentuais, e ainda operar frações.
\resposta{alternativa E — $\frac{8}{35}$.}
\begin{distratores}
\alt{A} 0 é a resposta de quem compara os \emph{valores} (que são os mesmos nos dois
gráficos) em vez das alturas das colunas.
\alt{B} $\frac{1}{2}$ não corresponde a nenhuma comparação coerente das alturas; é um
valor ``redondo'' que atrai quem estima no olho.
\alt{C} $\frac{1}{5}$ é só a razão do Gráfico 2; faltou subtraí-la da razão do Gráfico 1.
\alt{D} $\frac{2}{15} = \frac{1}{3} - \frac{1}{5}$: aparece para quem também ``corta'' o
Gráfico 1, lendo as alturas a partir da linha do 10 ($\frac{20}{60} = \frac{1}{3}$).
\end{distratores}
\end{resolucao}

%% =====================================================================
\secaoineditas
%% =====================================================================

\begin{questaoinedita}{1}{12}
Aplicativos de exercício físico registram a distância percorrida e o tempo de
cada treino. Em um domingo, o aplicativo de uma ciclista exibiu o resumo a seguir.

\begin{center}
\begin{tikzpicture}
  \node[draw=ebookLine,line width=.8pt,rounded corners=6pt,fill=ebookSoft,inner sep=8pt,
    font=\sffamily,align=center,text width=6.2cm] {%
    {\small\color{ebookMuted}TREINO DE DOMINGO \enspace\faBicycle}\\[4pt]
    \begin{tabular}{@{}c@{\hspace{18pt}}c@{}}
      {\Large\bfseries 30,0 km} & {\Large\bfseries 1 h 12 min}\\
      {\scriptsize\color{ebookMuted}distância} & {\scriptsize\color{ebookMuted}tempo em movimento}
    \end{tabular}};
\end{tikzpicture}
\end{center}

A velocidade média dessa ciclista nesse treino, em quilômetro por hora, foi de
\begin{alternativas*}[5]
  \item 0,42.
  \item 25,0.
  \item 26,8.
  \item 30,0.
  \item 36,0.
\end{alternativas*}
\end{questaoinedita}
\begin{resolucao}{B}
\passo{1. Convertendo o tempo.} A velocidade deve sair em km/h, então o tempo precisa
estar em horas: $12~\text{min} = \frac{12}{60}~\text{h} = 0{,}2~\text{h}$. Logo,
$1~\text{h}\,12~\text{min} = 1{,}2~\text{h}$.

\passo{2. Calculando.}
\[
v_m = \frac{30~\text{km}}{1{,}2~\text{h}} = 25~\text{km/h}.
\]
\resposta{alternativa B — 25,0~km/h.}
\begin{distratores}
\alt{A} 0,42 é $30 \div 72$: o aluno usou o tempo em minutos (72~min) e obteve uma
velocidade em km/\emph{min}, não em km/h.
\alt{C} 26,8 é $30 \div 1{,}12$: leu ``1~h~12~min'' como 1,12~h. Mas 12~min são $0{,}2$~h,
e não $0{,}12$~h.
\alt{D} 30,0 ignora os 12 minutos, como se a ciclista tivesse pedalado exatamente 1~hora.
\alt{E} 36,0 é $30 \cdot 1{,}2$: multiplicou a distância pelo tempo, em vez de dividir.
\end{distratores}
\end{resolucao}

\begin{questaoinedita}{2}{16}
Em uma loja de tintas, o tom ``verde-folha'' é obtido misturando tinta amarela e
tinta azul na razão de $3 : 2$, nessa ordem, em volume. Um pintor precisa de
18 litros de verde-folha para pintar a fachada de uma casa. Ele já tem guardados
5 litros de tinta azul, que sobraram de outro serviço, e vai comprar apenas o que
faltar de cada cor.

Quantos litros de tinta azul esse pintor ainda precisa comprar?
\begin{alternativas*}[5]
  \item 2,2
  \item 4,0
  \item 5,8
  \item 7,0
  \item 7,2
\end{alternativas*}
\end{questaoinedita}
\begin{resolucao}{A}
\passo{1. Método das partes.} A razão amarela : azul $= 3 : 2$ indica $3 + 2 = 5$ partes
iguais. Para 18~L, cada parte vale $18 \div 5 = 3{,}6$~L.

\passo{2. Tinta azul necessária.} Azul $= 2 \cdot 3{,}6 = 7{,}2$~L (e amarela
$= 3 \cdot 3{,}6 = 10{,}8$~L).

\passo{3. Descontando o estoque.} Como o pintor já tem 5~L de azul, precisa comprar
$7{,}2 - 5 = 2{,}2$~L.
\resposta{alternativa A — 2,2~litros.}
\begin{distratores}
\alt{B} 4,0 é $9 - 5$: o aluno dividiu os 18~L igualmente entre as duas cores (9~L de
cada), ignorando a razão $3 : 2$.
\alt{C} 5,8 é $10{,}8 - 5$: calculou corretamente as partes, mas usou a quantidade de
tinta \emph{amarela} (3 partes) no lugar da azul.
\alt{D} 7,0 é $12 - 5$: tratou a razão parte-parte $\frac{2}{3}$ como fração do todo,
calculando $\frac{2}{3}$ de 18~L. O correto é $\frac{2}{5}$ de 18~L.
\alt{E} 7,2~L é a quantidade total de azul da mistura; faltou descontar os 5~L que o
pintor já tinha.
\end{distratores}
\end{resolucao}

\begin{questaoinedita}{2}{16}
A bula de um antibiótico em suspensão oral informa que cada 5~mL do medicamento
contêm 250~mg do princípio ativo. Para uma criança com 24~kg de massa corporal, o
pediatra prescreveu 25~mg do princípio ativo por quilograma de massa corporal por
dia, divididos em doses iguais, administradas de 8 em 8 horas.

O volume, em mililitro, de cada dose que essa criança deve tomar é
\begin{alternativas*}[5]
  \item 0,8.
  \item 1,5.
  \item 4,0.
  \item 12,0.
  \item 40,0.
\end{alternativas*}
\end{questaoinedita}
\begin{resolucao}{C}
\passo{1. Dose diária.} $25~\tfrac{\text{mg}}{\text{kg}} \cdot 24~\text{kg} = 600$~mg por dia.

\passo{2. Número de doses.} De 8 em 8 horas são $24 \div 8 = 3$ doses por dia. Cada
dose tem $600 \div 3 = 200$~mg.

\passo{3. Concentração e volume.} A concentração é
$\frac{250~\text{mg}}{5~\text{mL}} = 50~\text{mg/mL}$. Logo, cada dose ocupa
\[
\frac{200~\text{mg}}{50~\text{mg/mL}} = 4~\text{mL}.
\]
\resposta{alternativa C — 4,0~mL.}
\begin{distratores}
\alt{A} 0,8 é $200 \div 250$: o aluno tratou a concentração como 250~mg \emph{por mL},
esquecendo que os 250~mg estão em 5~mL.
\alt{B} 1,5 vem de dividir a dose diária por 8 ($600 \div 8 = 75$~mg, que ocupam 1,5~mL):
confundiu ``de 8 em 8 horas'' com ``8 doses por dia''.
\alt{D} 12,0~mL é o volume da dose \emph{diária} ($600 \div 50$); faltou dividir em 3 doses.
\alt{E} 40,0 é $200 \div 5$: dividiu a dose em miligramas pelo volume de 5~mL, sem
construir a concentração corretamente.
\end{distratores}
\end{resolucao}

\begin{questaoinedita}{3}{17}
Um entregador sai do centro de distribuição de uma empresa e vai até uma cidade
vizinha, a 90~km de distância, dirigindo a uma velocidade média de 90~km/h. Lá,
ele permanece 30 minutos descarregando as encomendas e, em seguida, retorna ao
centro de distribuição pelo mesmo caminho, a uma velocidade média de 60~km/h,
por causa do trânsito intenso.

Considerando todo o período entre a saída e a chegada de volta ao centro de
distribuição, incluindo o tempo de descarga, a velocidade média do entregador,
em quilômetro por hora, foi de
\begin{alternativas*}[5]
  \item 50,0.
  \item 60,0.
  \item 64,3.
  \item 72,0.
  \item 75,0.
\end{alternativas*}
\end{questaoinedita}
\begin{resolucao}{B}
\passo{1. Tempo de cada etapa.} Ida: $\frac{90~\text{km}}{90~\text{km/h}} = 1$~h.
Descarga: 30~min $= 0{,}5$~h. Volta: $\frac{90~\text{km}}{60~\text{km/h}} = 1{,}5$~h.

\passo{2. Totais.} Distância: $90 + 90 = 180$~km. Tempo: $1 + 0{,}5 + 1{,}5 = 3$~h.

\passo{3. Velocidade média.}
\[
v_m = \frac{180~\text{km}}{3~\text{h}} = 60~\text{km/h}.
\]
A velocidade média é sempre distância total sobre tempo total — e o enunciado pediu
explicitamente para incluir a descarga.
\resposta{alternativa B — 60,0~km/h.}
\begin{distratores}
\alt{A} 50,0 é $\frac{90 + 0 + 60}{3}$: média aritmética das velocidades das três etapas,
contando a parada como ``velocidade zero''. As etapas têm durações diferentes, então
essa média não faz sentido.
\alt{C} 64,3 é $180 \div 2{,}8$: o aluno escreveu 30~min como 0,3~h.
\alt{D} 72,0 é $180 \div 2{,}5$: velocidade média correta \emph{sem} contar a descarga,
contrariando o enunciado.
\alt{E} 75,0 é $\frac{90 + 60}{2}$: a média aritmética das velocidades de ida e volta,
o erro mais comum. Mesmo sem a parada, ela não seria a velocidade média (que daria 72).
\end{distratores}
\end{resolucao}

\begin{questaoinedita}{3}{5}
Uma empresa de entregas vai comprar um veículo utilitário e escolherá o modelo com o
\textbf{menor custo de combustível por quilômetro rodado}. Os fabricantes, porém, não
informam o consumo da mesma maneira: alguns indicam quantos quilômetros o veículo
percorre com uma unidade de combustível; outros, quantos litros ele gasta para
percorrer 100 quilômetros. O quadro reúne essas informações e o preço do combustível
usado por cada modelo.

\begin{center}
\small
\begin{tabular}{cccc}
\toprule
\textbf{Modelo} & \textbf{Combustível} & \textbf{Consumo informado} & \textbf{Preço do combustível}\\
\midrule
I   & Gasolina & 15~km/L       & R\$~6,30 por litro\\
II  & Etanol   & 11~km/L       & R\$~4,10 por litro\\
III & Gasolina & 6,0~L/100\,km & R\$~6,30 por litro\\
IV  & Diesel   & 7,0~L/100\,km & R\$~5,20 por litro\\
V   & GNV      & 9~km/m³       & R\$~3,80 por metro cúbico\\
\bottomrule
\end{tabular}
\end{center}

O modelo escolhido pela empresa será o
\begin{alternativas*}[5]
  \item I.
  \item II.
  \item III.
  \item IV.
  \item V.
\end{alternativas*}
\end{questaoinedita}
\begin{resolucao}{D}
\passo{1. A taxa que importa.} Queremos reais por quilômetro (R\$/km). Para quem informa
km/L: custo $=\dfrac{\text{preço por litro}}{\text{km por litro}}$. Para quem informa
L/100\,km: custo $=\dfrac{\text{litros para 100 km} \times \text{preço}}{100}$.

\passo{2. Calculando.}
\begin{center}
\begin{tabular}{ll}
I: $6{,}30 \div 15 = 0{,}420$ & II: $4{,}10 \div 11 \approx 0{,}373$\\[2pt]
III: $6{,}0 \cdot 6{,}30 \div 100 = 0{,}378$ & IV: $7{,}0 \cdot 5{,}20 \div 100 = 0{,}364$\\[2pt]
V: $3{,}80 \div 9 \approx 0{,}422$ &
\end{tabular}
\end{center}

\passo{3. Conclusão.} O menor custo é o do modelo IV: cerca de R\$~0,36 por km.
Em \num{10000}~km, a diferença para o modelo II já passa de R\$~90.
\resposta{alternativa D — modelo IV.}
\begin{distratores}
\alt{A} O modelo I tem o maior rendimento entre os informados em km/L (15~km/L), mas usa
o combustível mais caro: custa R\$~0,42 por km. Rendimento sozinho não decide.
\alt{B} Quem lê ``6,0~L/100\,km'' e ``7,0~L/100\,km'' como se fossem km/L acha que III e IV
são péssimos ($6{,}30 \div 6 = 1{,}05$ e $5{,}20 \div 7 \approx 0{,}74$~R\$/km) e escolhe o
melhor entre os demais, o II ($\approx 0{,}373$).
\alt{C} O modelo III tem, de fato, o maior rendimento de todos ($100 \div 6 \approx 16{,}7$~km/L)
e o menor número na coluna de consumo, mas a gasolina cara o deixa em R\$~0,378 por km.
\alt{E} O modelo V usa o combustível de menor preço por unidade (R\$~3,80), porém rende
apenas 9~km por m³: R\$~0,42 por km, o mais caro.
\end{distratores}
\end{resolucao}

\begin{questaoinedita}{4}{18}
Para comparar a ocorrência de uma doença em locais com populações diferentes, os
órgãos de saúde usam a \textbf{taxa de incidência}, definida como o número de casos
novos registrados em um ano para cada 100 mil habitantes.

Uma região de saúde é formada apenas por dois municípios, P e Q. No último ano, o
município P, com 400 mil habitantes, teve taxa de incidência de 25 casos por 100 mil
habitantes, e o município Q, com 100 mil habitantes, teve taxa de 70 casos por
100 mil habitantes. Para o próximo ano, a meta é que a taxa de incidência da região
como um todo seja de, no máximo, 30 casos por 100 mil habitantes. Os gestores preveem
que as populações dos dois municípios e a taxa de incidência de P não se alterarão;
por isso, todas as ações de prevenção serão concentradas em Q.

Para que a meta seja cumprida, a taxa de incidência do município Q, no próximo ano,
deverá ser de, no máximo,
\begin{alternativas}
  \item 10 casos por 100 mil habitantes.
  \item 20 casos por 100 mil habitantes.
  \item 30 casos por 100 mil habitantes.
  \item 35 casos por 100 mil habitantes.
  \item 50 casos por 100 mil habitantes.
\end{alternativas}
\end{questaoinedita}
\begin{resolucao}{E}
\passo{1. Transformando taxas em casos.} P tem $400$ mil $= 4 \times 100$ mil habitantes;
com 25 casos por 100 mil, são $4 \cdot 25 = 100$ casos por ano. Q tem hoje
$1 \cdot 70 = 70$ casos. A região tem 500 mil habitantes e $170$ casos: taxa atual de
$170 \div 5 = 34$ casos por 100 mil.

\passo{2. Limite de casos da região.} Para a taxa da região ser no máximo 30 por 100 mil,
com 500 mil habitantes, o total de casos pode ser no máximo $5 \cdot 30 = 150$.

\passo{3. Limite para Q.} Como P continuará com 100 casos, Q pode ter no máximo
$150 - 100 = 50$ casos em seus 100 mil habitantes: taxa máxima de
\textbf{50 casos por 100 mil habitantes}.

\passo{4. Por que é difícil.} A taxa da região é uma média \emph{ponderada} pelas
populações, e não a média simples das taxas. Como P é quatro vezes mais populoso, sua
taxa baixa ``pesa'' muito, e Q não precisa chegar a 30.
\resposta{alternativa E — 50 casos por 100 mil habitantes.}
\begin{distratores}
\alt{A} 10 é $50 \div 5$: o aluno encontrou corretamente os 50 casos permitidos em Q, mas
calculou a taxa usando a população da \emph{região} (500 mil) em vez da de Q.
\alt{B} 20 é a \emph{redução} necessária (de 70 para 50 casos), e não a nova taxa.
\alt{C} 30 supõe que cada município, isoladamente, precisa atingir a meta da região.
\alt{D} 35 vem de tratar a taxa da região como média simples: $\frac{25 + x}{2} = 30$.
Isso ignora que as populações são diferentes.
\end{distratores}
\end{resolucao}
